\documentclass[12pt]{article}
\usepackage[utf8]{inputenc}
\usepackage[T1]{fontenc}
\usepackage{lmodern}
\usepackage{amsmath,amssymb,amsthm,mathtools}
\usepackage{geometry}
\usepackage{booktabs}
\usepackage{array}
\usepackage{longtable}
\usepackage{enumitem}
\usepackage{hyperref}
\geometry{margin=1in}
\setlist[itemize]{leftmargin=1.4em}
\setlist[enumerate]{leftmargin=1.6em}

\title{Cyclic Oriented Graph Counting\\\large Complete Guided Notes for Your Project}
\author{}
\date{March 5, 2026}

\theoremstyle{definition}
\newtheorem{definition}{Definition}[section]
\newtheorem{remark}[definition]{Remark}
\newtheorem{example}[definition]{Example}
\newtheorem{proposition}[definition]{Proposition}
\newtheorem{theorem}[definition]{Theorem}
\newtheorem{lemma}[definition]{Lemma}
\newtheorem{corollary}[definition]{Corollary}

\begin{document}
\maketitle

\begin{abstract}
This note explains the cyclic counting part of your project in a proof-style but beginner-friendly way. It defines every object, every symmetry, and every formula used by the code path \texttt{count\_adjacency\_burnside(n)} in \texttt{cyclic\_sequences.py}. It also includes a full worked example for \(n=8\), matching the exact computed values \(288\), \(64\), and final answer \(22\).
\end{abstract}

\tableofcontents

\section{What Problem Is Being Solved?}
We study directed cycles with \(n\) edges. Each edge gets one of two orientations:
\begin{itemize}
    \item bit \(0\): edge points from vertex \(i\) to \(i+1\) (mod \(n\));
    \item bit \(1\): edge points from vertex \(i+1\) to \(i\).
\end{itemize}
So each object is a binary sequence \((e_0,\dots,e_{n-1})\in\{0,1\}^n\).

The target is not the raw number \(2^n\). We count \emph{distinct objects up to symmetry} (rotations and reflections). Optionally, we enforce constraints such as adjacency and primitivity.

\section{Core Objects and Definitions}
\subsection{Cycle and indexing}
Vertices are \(V=\{0,1,\dots,n-1\}\), arranged cyclically. Edge \(i\) connects vertex \(i\) and vertex \(i+1\pmod n\).

\subsection{Vertex signature}
At each vertex \(v\), two incident edges determine a signature in \(\{-2,0,+2\}\):
\[
\sigma(v)=2\cdot(\text{number of incoming incident edges at }v)-2.
\]
So:
\begin{itemize}
    \item \(-2\): both incident edges point out of \(v\),
    \item \(0\): one in and one out,
    \item \(+2\): both point into \(v\).
\end{itemize}

\subsection{Adjacency constraint (project default fast path)}
The code enforces:
\[
\text{No consecutive nonzero signatures with the same sign.}
\]
So \((+2,+2)\) and \((-2,-2)\) adjacent are forbidden, while \((+2,-2)\), \((0,+2)\), etc. are allowed.

\section{Symmetry: Why We Need Group Actions}
\begin{definition}[Dihedral group \(D_n\)]
\(D_n\) is the symmetry group of an \(n\)-gon. It has:
\begin{itemize}
    \item \(n\) rotations,
    \item \(n\) reflections,
\end{itemize}
for total size \(|D_n|=2n\).
\end{definition}

\begin{definition}[Action on edge sequences]
Your code uses:
\begin{itemize}
    \item rotations: cyclic shifts of the bit sequence,
    \item reflections: reverse order \emph{and} flip bits \(0\leftrightarrow1\).
\end{itemize}
The bit flip is necessary because reflection reverses geometric direction of each edge.
\end{definition}

\begin{definition}[Equivalent objects]
Two sequences are equivalent if one is obtained from the other by some element of \(D_n\).
\end{definition}

\begin{definition}[Orbit]
The orbit of an object is all sequences reachable by these symmetries. One orbit = one distinct graph in your count.
\end{definition}

\section{Canonical Representative}
For brute checks, the program computes a canonical form:
\begin{enumerate}
    \item generate all rotations of the sequence,
    \item generate all rotations of reflected-and-bit-flipped sequence,
    \item choose lexicographically smallest candidate.
\end{enumerate}
This gives one deterministic representative per orbit.

\section{Burnside Lemma: The Orbit Counter}
\subsection{Statement}
\begin{theorem}[Burnside]
For a finite group \(G\) acting on finite set \(X\),
\[
\#(X/G)=\frac{1}{|G|}\sum_{g\in G}|\mathrm{Fix}(g)|,
\]
where \(\mathrm{Fix}(g)=\{x\in X:\ g\cdot x=x\}\).
\end{theorem}

\subsection{What ``fixed'' means}
For one symmetry \(g\), an object is \emph{fixed} if applying \(g\) leaves it exactly unchanged.

Example: if rotating by 2 steps maps a sequence to itself, that sequence contributes to \(|\mathrm{Fix}(g)|\) for that rotation.

\subsection{Why this is correct (double-counting sketch)}
Count pairs \((g,x)\) with \(g\cdot x=x\) in two ways:
\begin{itemize}
    \item by symmetry first: \(\sum_g |\mathrm{Fix}(g)|\),
    \item by object first: \(\sum_x |\mathrm{Stab}(x)|\).
\end{itemize}
Using orbit-stabilizer, \(|\mathrm{Stab}(x)|\) summed over each orbit equals \(|G|\), so total is \(|G|\times\#\text{orbits}\). Rearranging yields Burnside.

\section{Why Burnside Alone Is Not Enough}
Burnside tells \emph{how to combine} fixed counts. It does not magically compute those counts.

So we still need fast methods for \(|\mathrm{Fix}(g)|\). That is where transfer matrices enter.

\section{Transfer Matrix for Adjacency-valid Cycles}
\subsection{State space used in your code}
A local state is:
\[
s=(v,e),\quad v\in\{-2,0,+2\},\ e\in\{0,1\}.
\]
So there are 6 states total.

Interpretation:
\begin{itemize}
    \item \(e\) is current edge bit,
    \item \(v\) is signature of the previous vertex.
\end{itemize}

\subsection{One-step update}
Given current edge \(e\) and next edge \(e'\), the next vertex signature is
\[
v' = 2\big(\mathbf{1}_{e=0}+\mathbf{1}_{e'=1}\big)-2.
\]
Transition is allowed only if adjacency rule holds between \(v\) and \(v'\).

\subsection{Matrix construction}
Define \(M\in\mathbb{Z}_{\ge 0}^{6\times 6}\) by
\[
M_{ij}=\begin{cases}
1,&\text{if transition state }i\to j\text{ is allowed},\\
0,&\text{otherwise.}
\end{cases}
\]
In your state order
\[
(-2,0),(-2,1),(0,0),(0,1),(2,0),(2,1),
\]
this matrix is
\[
M=
\begin{pmatrix}
0&0&1&0&0&1\\
0&0&0&1&0&0\\
0&0&1&0&0&1\\
1&0&0&1&0&0\\
0&0&1&0&0&0\\
1&0&0&1&0&0
\end{pmatrix}.
\]

\subsection{Closed walks and trace}
\begin{lemma}
\((M^m)_{ij}\) equals the number of length-\(m\) valid state paths from \(i\) to \(j\).
\end{lemma}

\begin{corollary}
The number of valid length-\(m\) closed state walks is
\[
C(m)=\mathrm{tr}(M^m).
\]
\end{corollary}

This \(C(m)\) is the key primitive used in Burnside rotation counts.

\section{Burnside Rotation Term in This Project}
For a rotation by \(k\), fixed sequences have period \(d=\gcd(n,k)\). Grouping rotations by \(d\) yields:
\[
\text{RotSum}(n)=\sum_{d\mid n} \varphi\!\left(\frac{n}{d}\right) C(d),
\]
where \(\varphi\) is Euler totient.

Interpretation:
\begin{itemize}
    \item \(C(d)\): how many base periodic blocks of length \(d\) satisfy local rules,
    \item \(\varphi(n/d)\): how many rotations have gcd class corresponding to that \(d\).
\end{itemize}

\section{Burnside Reflection Term in This Project}
Your code has a separate routine for reflections:
\begin{itemize}
    \item if \(n\) is odd: reflection-fixed count is 0,
    \item if \(n\) is even: compute half-length transfer and enforce reflection boundary compatibility.
\end{itemize}
Define this code-computed number as \(R(n)\).

Then reflection contribution to Burnside sum is:
\[
\text{RefTerm}(n)=\frac{n}{2}\,R(n),
\]
matching your implementation.

\section{Final Counting Formula in Code}
For \(n\ge1\):
\[
a(n)=\frac{\text{RotSum}(n)+\text{RefTerm}(n)}{2n}
=\frac{1}{2n}\left(\sum_{d\mid n}\varphi\!\left(\frac{n}{d}\right)\mathrm{tr}(M^d)+\frac{n}{2}R(n)\right).
\]
This is exactly the fast adjacency count route in your cyclic code.

\section{Worked Example: \(n=8\) (Exact Values From Your Program)}
Divisors of 8 are \(1,2,4,8\).

\begin{center}
\begin{tabular}{cccc}
\toprule
\(d\) & \(\varphi(8/d)\) & \(C(d)=\mathrm{tr}(M^d)\) & contribution \\
\midrule
1 & 4 & 2   & 8   \\
2 & 2 & 4   & 8   \\
4 & 1 & 16  & 16  \\
8 & 1 & 256 & 256 \\
\bottomrule
\end{tabular}
\end{center}

So
\[
\text{RotSum}(8)=8+8+16+256=288.
\]

From reflection routine:
\[
R(8)=16,
\quad
\text{RefTerm}(8)=\frac{8}{2}\cdot16=64.
\]

Final Burnside average:
\[
a(8)=\frac{288+64}{16}=\frac{352}{16}=22.
\]
So your program returns \(22\) distinct oriented cyclic graphs (with adjacency rule).

\section{Optional Constraints in the Cyclic Module}
\subsection{Primitivity}
A sequence is primitive if it is not a repetition of a shorter block.

Code fast formula uses Mobius inversion:
\[
a_{\mathrm{prim}}(n)=\sum_{d\mid n}\mu(d)\,a\!\left(\frac{n}{d}\right),
\]
where \(\mu\) is the Mobius function.

\subsection{Forbidden subsequence mode}
A signature is rejected if any cyclic contiguous block matches a canonical forbidden pattern learned from smaller lengths. This route uses explicit set-based enumeration (not the adjacency fast matrix shortcut).

\section{Validation and Correctness Strategy}
Your pipeline validates in layers:
\begin{enumerate}
    \item Brute-force canonical enumeration for small \(n\).
    \item Fast Burnside+matrix counts for adjacency mode.
    \item Assert equality on overlapping range.
    \item Compare unconstrained small terms with OEIS A053656 values.
\end{enumerate}

This is strong because independent methods agree on the same outputs.

\section{Complexity Intuition}
Brute force is \(O(2^n)\), so it explodes quickly.

Fast path uses:
\begin{itemize}
    \item divisor loop (depends on number of divisors),
    \item totient and Mobius on divisor arguments,
    \item matrix powers (small fixed size matrix, logarithmic exponentiation depth).
\end{itemize}
Hence very large \(n\) is feasible.

\section{Presentation-ready 90-second summary}
``I model each directed cycle as a binary edge sequence and quotient by dihedral symmetry, so I count true isomorphism classes, not raw strings. Burnside lemma says the number of classes is the average fixed-point count over all symmetries. To compute those fixed counts efficiently, I use a transfer matrix on 6 local states that encode the previous vertex signature and current edge orientation, with adjacency constraint built into allowed transitions. Rotation-fixed counts reduce to weighted trace terms over divisors, and reflection-fixed counts are computed by half-length transfer with boundary compatibility. For example at \(n=8\), rotation sum is 288 and reflection term is 64, giving \((288+64)/16=22\) distinct classes. Small \(n\) is cross-checked by brute-force canonical enumeration and OEIS values.''

\section{Glossary of Symbols}
\begin{longtable}{>{\raggedright\arraybackslash}p{0.22\linewidth}p{0.72\linewidth}}
\toprule
Symbol & Meaning \\
\midrule
\(n\) & cycle length (number of edges and vertices) \\
\(e_i\) & edge orientation bit at edge \(i\), in \(\{0,1\}\) \\
\(\sigma(v)\) & vertex signature at vertex \(v\), in \(\{-2,0,+2\}\) \\
\(D_n\) & dihedral group of order \(2n\) \\
\(G\) & symmetry group acting on objects \\
\(\mathrm{Fix}(g)\) & objects unchanged by symmetry \(g\) \\
\(M\) & transfer matrix on 6 local states \\
\(C(m)\) & \(\mathrm{tr}(M^m)\), number of length-\(m\) closed state walks \\
\(\varphi\) & Euler totient function \\
\(\mu\) & Mobius function \\
\(R(n)\) & reflection-fixed count computed by reflection routine \\
\(a(n)\) & final orbit count (distinct cyclic graphs) \\
\bottomrule
\end{longtable}

\section{Takeaway}
Your cyclic module is mathematically clean:
\begin{enumerate}
    \item define equivalence by dihedral action,
    \item compute fixed counts efficiently by transfer matrices,
    \item assemble exact orbit counts by Burnside,
    \item validate against brute force and known sequences.
\end{enumerate}
That is why the method is both rigorous and scalable.

\end{document}
